\section{Marco Teórico}

La obsolescencia programada está definida como “el uso de técnicas mediante las
cuales una entidad responsable de lanzar un producto en el mercado busca reducir
el ciclo de vida útil —de manera deliberada— de dicho producto”
\parencite[p.~741]{malinauskaite2021}. Esta se manifiesta de distintas formas.
Una de ellas consiste en diseñar productos difíciles o imposibles de reparar, ya
sea porque las piezas necesarias no están disponibles en el mercado o porque el
costo de la reparación resulta tan elevado que adquirir uno nuevo termina siendo
la opción más práctica —no porque la reparación sea técnicamente imposible, sino
porque el fabricante decidió restringirla mediante limitaciones digitales
\parencite[p.~105]{alfieri2023}.

Según \textcite[p.~225]{kirchherr2017}, la economía circular “describe un
sistema económico basado en modelos de negocio que sustituyen el concepto de
``fin de vida útil'' con el propósito de alcanzar un desarrollo sustentable”. En
la práctica, este modelo se sostiene sobre nueve principios fundamentales
—rechazar, replantear, reducir, reutilizar, reparar, restaurar, reciclar,
recuperar, reusar— orientados a generar calidad ambiental y prosperidad
económica.

Entre los principios de la economía circular, la reparación, restauración, así
como el replanteamiento resultan particularmente relevantes como estrategias
para contrarrestar la lógica de reemplazo acelerado que sustenta la
obsolescencia programada. Estos principios, según
\textcite[pp.~29-31]{lambiase2024} están definidos de la siguiente forma:  

\begin{enumerate}
    \item La \textbf{reparabilidad} de productos es un principio orientado a
    extender su vida útil mediante el mantenimiento. Su implementación en el
    mercado electrónico puede tomar dos vías: la intervención de talleres de
    reparación especializados con acceso a las piezas necesarias, o la
    reparación por parte del propio usuario, siempre que la información técnica
    y las herramientas requeridas sean accesibles. 
    \item La \textbf{restauración} implica reemplazar componentes viejos o
    desgastados por partes nuevas, extendiendo así la vida útil del producto.
    Este principio comparte con la reparabilidad una condición para su
    implementación exitosa: que las herramientas, piezas necesarias sean
    accesibles para quienes buscan llevarla a cabo. Sin embargo, también depende
    de decisiones tomadas desde la etapa de diseño —los productos deben
    diseñarse bajo un esquema modular que permita retirar, reparar, reemplazar
    sus componentes de forma sencilla. 
    \item FInalmente, el \textbf{replanteamiento del consumo} se materializa en
    el modelo ``product as a service'' —producto como servicio—, en el que la
    empresa retiene la propiedad del bien, ofreciendo al cliente únicamente su
    uso mediante arrendamiento, renta o pago por uso. Bajo este esquema, los
    fabricantes tienen un incentivo directo para diseñar productos más
    duraderos, lo que reduce el desperdicio asociado al reemplazo constante. 
\end{enumerate}

Sin embargo, el modelo de producción dominante actual opera en sentido
contrario, priorizando el consumo a corto plazo y llevando al planeta hacia una
situación cada vez más insostenible. Este ritmo de producción acelerado
—activamente sostenido por prácticas como la obsolescencia programada— resulta
incompatible con los fundamentos de la economía circular, lo que convierte su
regulación en una condición necesaria para alcanzar los objetivos que este
modelo propone. 

Otra variante es la obsolescencia tecnológica o de incompatibilidad, que ocurre
cuando la introducción de un producto técnicamente más avanzado vuelve obsoletos
los modelos anteriores —por ejemplo, mediante actualizaciones de software
incompatibles o tan voluminosas que obligan al usuario a adquirir una versión
más reciente del dispositivo \parencite[p.~724]{malinauskaite2021}. El usuario
se encuentra entonces con un dispositivo que físicamente funciona, pero que no
puede utilizarse para las tareas básicas de la vida diaria. La presión por
mantenerse conectado termina empujándolo a adquirir un modelo más reciente,
aunque el anterior aún sea funcional. 

Mientras que la obsolescencia programada es un fenómeno ocasionado
principalmente por la industria, los consumidores también juegan un papel
importante, puesto que la obsolescencia percibida fomenta el reemplazo
innecesario de productos completamente funcionales
\parencite[pp.~6-7]{corrocher2025}. Sin embargo, atribuir este comportamiento
únicamente a una decisión libre del consumidor sería simplista. La economía
conductual señala que las personas suelen tomar decisiones irracionales debido a
su ``racionalidad limitada'', lo que las hace susceptibles a ser manipuladas en
ciertos contextos \parencite[pp.~725-726]{malinauskaite2021}.

En el mercado de productos electrónicos, esto se traduce en una sobrecarga
cognitiva provocada por la cantidad, así como complejidad de la información
disponible, dificultando que el consumidor pueda comparar y evaluar sus opciones
de forma efectiva —un fenómeno conocido como ``confusopolio''. En este fenómeno,
el consumidor se ve saturado de información deliberadamente confusa que eleva
los costos cognitivos de búsqueda, dificultando su capacidad de tomar decisiones
de compra verdaderamente informadas. 

Frente a esta problemática de sobrecarga informativa, resulta importante
examinar las herramientas que actualmente existen para facilitar el acceso a
información sobre productos electrónicos. BARCODE LOOKUP representa una
aproximación a este problema mediante el escaneo del código de barras de
cualquier producto. Su base de datos internacional —que cuenta con más de mil
millones de productos registrados provenientes— permite al usuario obtener, de
forma inmediata, información como nombre del producto, fabricante, descripción e
imágenes. Adicionalmente, la plataforma vincula al usuario con los vendedores en
línea que comercializan el mismo artículo, facilitando la comparación de precios
entre distintas tiendas. Sus datos provienen de una red de más de 1,500
minoristas a nivel mundial, así como de contribuciones directas de usuarios, lo
que le permite mantenerse actualizada de forma continua. 

Otra herramienta orientada a reducir la sobrecarga informativa del consumidor es
ShopSavvy \parencite{shopsavvy2008}, una aplicación móvil disponible para iOS y
Android que permite escanear el código de barras de cualquier producto para
acceder a precios en tiendas locales. Sin embargo, su aportación más relevante
en el contexto del confusopolio es otra: la integración de opiniones de usuarios
reales. A diferencia de las especificaciones técnicas —cuya comprensión exige un
nivel de conocimiento que no todos los consumidores poseen—, las reseñas ofrecen
información basada en experiencia directa de uso, lo que permite evaluar el
rendimiento real de un producto sin necesidad de entender terminología
especializada. En un mercado donde la complejidad informativa es, en muchos
casos, deliberada, este tipo de evaluación colectiva representa una vía
accesible para que el consumidor tome decisiones más informadas y,
potencialmente, menos susceptibles a la manipulación que sustenta la
obsolescencia percibida. 

Asimismo, el greenwashing representa otra fuente de desinformación que
profundiza la confusión del consumidor. \textcite[p.~196]{font2025} lo define
como el uso indebido de la comunicación sustentable para engañar al consumidor
sobre las características sostenibles de un producto. Esta práctica resulta
especialmente problemática porque explota las expectativas legítimas del
consumidor, quien asume que una empresa con discurso ecológico implementa
programas ambientales reales, adopta iniciativas sociales, opera bajo criterios
de sostenibilidad económica. Cuando estas expectativas no corresponden a la
realidad, el resultado es un escepticismo generalizado que le dificulta al
consumidor distinguir entre marcas genuinamente responsables y aquellas que
simplemente lo aparentan. 

Ante este panorama de desconfianza generado por el greenwashing, GS1 representa
una fuente de datos que opera en sentido contrario: en lugar de comunicar lo que
las empresas quieren mostrar, gestiona información verificable sobre productos.
GS1 es una organización global sin fines de lucro presente en más de 150 países
que desarrolla y administra estándares internacionales de identificación de
productos —siendo el código de barras su herramienta más reconocida. A través de
su sistema de identificación único por producto, conocido como Número Global de
Artículo Comercial (GTIN), es posible acceder a datos sobre el origen de los
materiales, el proceso de fabricación las certificaciones de sostenibilidad
reales asociadas a un producto y su fabricante. Esta última función resulta
especialmente significativa: GS1 Sweden, en colaboración con TCO Certified
—organismo líder en certificación de sustentabilidad para productos de
tecnología— ha integrado el GTIN en sus procesos de certificación, permitiendo
que productos electrónicos desarrollados puedan ser rastreados y vinculados
directamente a sus certificaciones verificadas \parencite{gs12025}. Esto
convierte a GS1 en una herramienta con potencial para contrarrestar el
greenwashing, al ofrecer al consumidor un medio de verificación objetivo frente
a los discursos corporativos que no siempre corresponden a la realidad. 

En un contexto comercial donde este tipo de problemáticas influyen negativamente
en las decisiones de compra, toma especial importancia la existencia de fuentes
de información no solo confiables, sino presentadas de forma comprensible para
cualquier tipo de consumidor, independientemente de su conocimiento previo sobre
tecnología. La teoría de la riqueza de los medios sugiere que la cantidad
adecuada de información en un estímulo puede reducir la incertidumbre e influir
positivamente en el resultado de la comunicación \parencite[p.~197]{font2025}.
Sin embargo, este equilibrio es delicado: si bien una mayor riqueza informativa
transmite más datos al receptor, un exceso puede derivar en sobrecarga
informativa. Esta ocurre cuando los requerimientos de procesamiento de
información superan la capacidad del individuo para procesarla, convirtiéndose
en un obstáculo para la toma de decisiones incluso cuando dicha información es
relevante \parencite[p.~197]{font2025}.  

La situación se agrava cuando la información es poco familiar, compleja o
presentada a un ritmo mayor al que el receptor puede asimilarla —condiciones que
caracterizan precisamente la forma en que los fabricantes de dispositivos
electrónicos comunican sus productos al mercado. 

Frente a la problemática de sobrecarga informativa previamente descrita, la
divulgación progresiva —progressive disclosure— emerge como una estrategia de
diseño de interfaz orientada a mitigar sus efectos en la toma de decisiones del
consumidor. Este enfoque consiste en presentar inicialmente un conjunto reducido
de atributos clave, ofreciendo acceso opcional a información más detallada en
etapas posteriores. \textcite[pp.~15-19]{hu2022} demostraron empíricamente que
los usuarios expuestos primero a conjuntos simplificados de información logran
establecer criterios de decisión más estables, los cuales mantienen al
enfrentarse posteriormente a datos más complejos. En contraste, quienes reciben
toda la información de forma simultánea exhiben una mayor inconsistencia entre
sus preferencias declaradas y sus decisiones reales —fenómeno atribuido
directamente al exceso de carga cognitiva. Aplicado al contexto de dispositivos
electrónicos, este principio sugiere que una aplicación que estructure la
información técnica y ambiental de manera progresiva no solo reduciría la
confusión del consumidor, sino que favorecería decisiones de compra más
alineadas con sus valores y necesidades reales.

La aplicación diseñada toma como eje central el cuidado de la carga cognitiva
del usuario, adoptando la divulgación progresiva como estrategia
estructural de presentación de información, permitiendo que el consumidor
procese los datos de manera gradual, construya criterios de decisión sólidos
antes de enfrentarse a información de mayor complejidad. Así pues, se busca
que las decisiones de compra sean producto de una comprensión real, no de una
saturación informativa que, como se ha descrito, puede derivar en elecciones
incongruentes con las necesidades y valores del propio usuario. 

Además, la aplicación reconoce el valor de la experiencia colectiva como fuente de
información accesible. A diferencia de las especificaciones técnicas
—cuya interpretación exige conocimientos especializados—, las opiniones de
usuarios reales ofrecen una evaluación basada en el uso cotidiano del producto.
La plataforma integró este tipo de reseñas bajo un modelo libre de
patrocinios, con el fin de preservar su credibilidad, contrarrestando la
distorsión informativa que caracteriza los canales comerciales convencionales.
Finalmente, retomando la problemática de la obsolescencia programada, la aplicación
incorporó un directorio de opciones de reparación cercanas al usuario. Esta
funcionalidad busca traducir los principios de reparabilidad y restauración de
la economía circular —previamente descritos— en una herramienta que
fomente su adopción, contribuyendo así a la concientización del consumidor y a
la extensión de la vida útil de los dispositivos electrónicos.